
Die Mikrobefehle der Registermaschine Johnny sind:

\begin{verbatim}
    db->ram  : Inhalt des Datenbus im RAM speichern
    ram->db  : Inhalt des RAM auf den Datenbus legen
    db->ins  : Inhalt des Datenbus ins Befehlsregister übernehmen
    db->acc  : Inhalt des Datenbus in den Akkumulator übernehmen
    acc->db  : Inhalt des Akkumulators auf den Datenbus legen
    plus     : Akkumulator um Inhalt des Datenbus erhöhen
    minus    : Akkumulator um Inhalt des Datenbus erniedrigen
    acc++    : Akkumulator um 1 erhöhen
    acc--    : Akkumulator um 1 erniedrigen
    acc=0    : Akkumulator auf 0 setzen
    pc->ab   : Befehlszähler auf den Adressbus legen
    pc++     : Befehlszähler um 1 erhöhen
    =0:pc++  : falls Akkumulator = 0, Befehlszähler um 1 erhöhen
    ins->ab  : Adressteil des Befehlsregisters auf den Adressbus legen
    ins->pc  : Adressteil des Befehlsregisters in den Befehlszähler schreiben
    ins->mc  : zum Mikroprogramm des Befehls im Befehlsregister springen
    mc=0     : zurück zum Anfang des Mikrocodes (FETCH)
    stop     : Programm beenden
\end{verbatim}

